\chapter{Ngân hàng bài tập có lời giải từng bước}\label{ch:exercises}
Các bài dưới đây là bài luyện tự biên soạn hoặc diễn giải từ dạng bài
trong các nguồn ở thư mục môn học, không được gán số bài của sách hay
đề thi chính thức. Dùng cùng quy ước kernel và quantile trong giáo trình.
Làm tròn ở bước cuối; giữ phân số hoặc biểu thức chuẩn tắc càng lâu càng tốt.

\section{Dạng 1: tính KDE và đạo hàm}
\begin{exercise}[Mẫu gốc và điểm không khả vi]\label{ex:original}
Cho $X=(1,2,4,5)$, $x=3$, $h=2$. Tính $\widehat f_2(3)$,
$\widehat f_2'(3)$, $\widehat f_2''(3)$ với Epanechnikov, Uniform,
Triangular và Gaussian. Giải thích nếu đạo hàm không tồn tại.
\end{exercise}
\begin{solution}
\step{1. Chuẩn hóa khoảng cách và các mẫu số.}
$n=4$, $(u_1,u_2,u_3,u_4)=(1,1/2,-1/2,-1)$.
Các mẫu số cho mật độ, đạo hàm một và hai là
$nh=8$, $nh^2=16$, $nh^3=32$.
\begin{center}
\begin{tabular}{ccccc}
\toprule $u$ & Epanechnikov & Uniform & Triangular & Gaussian\\\midrule
$1$ & $0$ & $1/2$ & $0$ & $\varphi(1)$\\
$1/2$ & $9/16$ & $1/2$ & $1/2$ & $\varphi(1/2)$\\
$-1/2$ & $9/16$ & $1/2$ & $1/2$ & $\varphi(1/2)$\\
$-1$ & $0$ & $1/2$ & $0$ & $\varphi(1)$\\\bottomrule
\end{tabular}
\end{center}
\step{2. Cộng trọng số và chia cho $nh$.}
\[
\widehat f_{\rm E}(3)=\frac{18/16}{8}=\boxed{\frac9{64}},\quad
\widehat f_{\rm U}(3)=\frac{2}{8}=\boxed{\frac14},\quad
\widehat f_{\rm T}(3)=\frac1{8}=\boxed{\frac18}.
\]
Với Gaussian,
\[
\widehat f_{\rm G}(3)=\boxed{\frac{\varphi(1)+\varphi(1/2)}4}
\approx0.14850901.
\]
Uniform dùng đầu mút đóng; nếu dùng $|u|<1$, giá trị tại đúng 3 sẽ
là $1/8$. Giá trị giới hạn quanh 3 vẫn phải kiểm tra riêng.

\step{3. Gaussian: dùng trực tiếp công thức đạo hàm.}
$K'$ lẻ, nên các cặp đối xứng triệt tiêu và $\widehat f_{\rm G}'(3)=0$.
$K''(\pm1)=0$, $K''(\pm1/2)=-3\varphi(1/2)/4$, nên
\[
\widehat f_{\rm G}''(3)=\frac{-3\varphi(1/2)/2}{32}
=\boxed{-\frac{3\varphi(1/2)}{64}}\approx-0.01650306.
\]

\step{4. Epanechnikov: xét tổng ở $3+\delta$.}
Với $|\delta|$ đủ nhỏ, hai điểm 2 và 4 ở trong cửa sổ. Tổng kernel
của chúng bằng $9/8-3\delta^2/8$.
Nếu $\delta>0$, chỉ điểm 5 ở đầu ngoài đi vào cửa sổ, đóng góp
$3\delta/4-3\delta^2/16$; nếu $\delta<0$, điểm 1 đi vào, cho
$-3\delta/4-3\delta^2/16$. Ghép hai phía:
\[
\widehat f_{\rm E}(3+\delta)=\frac9{64}
+\frac3{32}|\delta|-\frac9{128}\delta^2.
\]
Thương sai phân từ trái tiến về $-3/32$, từ phải về $3/32$.
Vì hai giới hạn khác nhau,
$\boxed{\widehat f_{\rm E}'(3)\text{ không tồn tại}}$;
đạo hàm bậc hai cổ điển cũng không tồn tại.
Tính đối xứng của hàm không buộc đạo hàm tại tâm tồn tại: $|\delta|$
là phản ví dụ đơn giản.

\step{5. Triangular và Uniform: kiểm tra tương tự.}
Với Triangular, hai điểm nội cho tổng 1; một điểm ở ranh giới đi vào cho
$|\delta|/2$:
\[
\widehat f_{\rm T}(3+\delta)=\frac18+\frac1{16}|\delta|.
\]
Hai độ dốc là $-1/16$ và $1/16$, nên cả hai đạo hàm yêu cầu không tồn tại.
Với Uniform đóng, tại 3 có bốn điểm được đếm, nhưng với
$0<|\delta|$ đủ nhỏ chỉ còn ba điểm. Vì thế
\[
\lim_{\delta\to0,\,\delta\ne0}\widehat f_{\rm U}(3+\delta)=\frac3{16}
\ne\frac14=\widehat f_{\rm U}(3).
\]
Hàm không liên tục tại 3, nên không có đạo hàm một hoặc hai ở đó.
\pitfall{Không dùng giá trị $K''=-3/2$ của Epanechnikov ở điểm nối;
không bỏ các điểm $u=\pm1$ khi kiểm tra giới hạn. Giá trị mật độ tại
đầu mút và tính khả vi của tổng là hai câu hỏi khác nhau.}
\end{solution}

\begin{exercise}[Cùng mẫu, điểm khả vi]
Giữ $X=(1,2,4,5)$ và $x=3$, đổi $h=3/2$. Tính ba đại lượng trên
cho cùng bốn kernel.
\end{exercise}
\begin{solution}
\step{1. Xác định điểm nội và ngoại.}
$u=(4/3,2/3,-2/3,-4/3)$. Với các kernel compact, hai điểm ngoài
không đóng góp, hai điểm còn lại không ở điểm nối hoặc ở 0.
Các mẫu số là $nh=6$, $nh^2=9$, $nh^3=27/2$.

\step{2. Epanechnikov.}
$K(\pm2/3)=\frac34(1-4/9)=5/12$, nên
\[
\widehat f(3)=\frac{2(5/12)}6=\boxed{\frac5{36}},\quad
\widehat f'(3)=\frac{-1+1}{9}=\boxed0,\quad
\widehat f''(3)=\frac{-3}{27/2}=\boxed{-\frac29}.
\]

\step{3. Uniform và Triangular.}
Uniform cho $\widehat f(3)=2(1/2)/6=\boxed{1/6}$,
hai đạo hàm bằng 0 vì estimator hằng trong một lân cận của 3.
Triangular cho $\widehat f(3)=2(1/3)/6=\boxed{1/9}$,
đạo hàm một bằng $(-1+1)/9=0$, đạo hàm hai bằng 0 tại điểm này.

\step{4. Gaussian.}
Đặt $a=\varphi(4/3)$, $b=\varphi(2/3)$:
\[
\widehat f(3)=\boxed{\frac{a+b}{3}}\approx0.16115269,\quad
\widehat f'(3)=\boxed0,
\]
\[
\widehat f''(3)=\frac{2[(7/9)a-(5/9)b]}{27/2}
=\boxed{\frac{4(7a-5b)}{243}}\approx-0.00739374.
\]
Ở đây áp dụng công thức đạo hàm hợp lệ vì tất cả đối số tránh điểm
không khả vi. Đây không phải chứng cứ rằng Uniform ước lượng $f''$ tốt:
đạo hàm bằng 0 chỉ phản ánh hàm hằng từng đoạn.
\end{solution}

\begin{exercise}[Biweight ở một điểm nội]
Cho $X=(0,1,2)$, $x=1$, $h=2$, kernel Biweight.
Tính mật độ, đạo hàm một và hai.
\end{exercise}
\begin{solution}
\step{1. Lập bảng ba giá trị.}
$u=(1/2,0,-1/2)$, mẫu số $(nh,nh^2,nh^3)=(6,12,24)$.
\[
K(\pm1/2)=\frac{15}{16}\left(\frac34\right)^2=\frac{135}{256},
\qquad K(0)=\frac{15}{16}.
\]
\step{2. Mật độ.}
\[
\widehat f_2(1)=\frac{270/256+240/256}{6}
=\boxed{\frac{85}{256}}.
\]
\step{3. Đạo hàm.}
$K'$ lẻ nên tổng bằng 0. Với $K''(u)=15(3u^2-1)/4$,
\[
K''(\pm1/2)=-\frac{15}{16},\qquad K''(0)=-\frac{15}{4},
\]
\[
\boxed{\widehat f_2'(1)=0,\qquad
\widehat f_2''(1)=\frac{-15/16-15/4-15/16}{24}
=-\frac{15}{64}.}
\]
Độ cong âm thể hiện hình dạng lõm xuống của estimator tại tâm mẫu này.
\end{solution}

\section{Dạng 2: tích phân hằng số kernel}
\begin{exercise}[Uniform và Triangular]
Tính trực tiếp $\mu_2(K)$ và $R(K)$ của hai kernel.
\end{exercise}
\begin{solution}
\step{1. Uniform.}
Đối xứng cho
\[
\mu_2=2\int_0^1u^2\frac12du=[u^3/3]_0^1=\boxed{1/3},
\quad
R=2\int_0^1(1/2)^2du=\boxed{1/2}.
\]
\step{2. Triangular.}
Trên $[0,1]$, $K=1-u$, nên
\[
\mu_2=2\int_0^1(u^2-u^3)du
=2\left(\frac13-\frac14\right)=\boxed{1/6},
\]
\[
R=2\int_0^1(1-2u+u^2)du
=2\left(1-1+\frac13\right)=\boxed{2/3}.
\]
Không bỏ hệ số 2 khi chuyển tích phân kernel chẵn từ $[-1,1]$ sang $[0,1]$.
\end{solution}

\begin{exercise}[Epanechnikov]
Kiểm tra chuẩn hóa, rồi tính $\mu_2$ và $R$.
\end{exercise}
\begin{solution}
\step{1. Diện tích.}
$\int K=(3/2)\int_0^1(1-u^2)du=(3/2)(1-1/3)=1$.
\step{2. Moment bậc hai.}
\[
\mu_2=\frac32\int_0^1(u^2-u^4)du
=\frac32\left(\frac13-\frac15\right)=\boxed{\frac15}.
\]
\step{3. Bình phương kernel.}
Phải bình phương cả hệ số $3/4$:
\[
R=\frac98\int_0^1(1-2u^2+u^4)du
=\frac98\left(1-\frac23+\frac15\right)
=\frac98\frac8{15}=\boxed{\frac35}.
\]
Hai hằng số này đi vào AMISE theo cách khác nhau: $\mu_2$ bình phương
trong squared bias, $R$ tuyến tính trong variance.
\end{solution}

\begin{exercise}[Biweight]
Tính $\mu_2$ và $R$ bằng khai triển đa thức.
\end{exercise}
\begin{solution}
\step{1. Chuẩn hóa.}
\[
\int K=\frac{15}{8}\int_0^1(1-2u^2+u^4)du
=\frac{15}{8}\frac8{15}=1.
\]
\step{2. Moment bậc hai.}
\[
\mu_2=\frac{15}{8}\int_0^1(u^2-2u^4+u^6)du
=\frac{15}{8}\left(\frac13-\frac25+\frac17\right)
=\frac{15}{8}\frac8{105}=\boxed{\frac17}.
\]
\step{3. Bình phương.}
$(1-u^2)^4=1-4u^2+6u^4-4u^6+u^8$, nên
\begin{align*}
R&=\frac{225}{128}\left(1-\frac43+\frac65-\frac47+\frac19\right)\\
&=\frac{225}{128}\frac{128}{315}=\boxed{\frac57}.
\end{align*}
Giữ phân số tránh sai số khi các số hạng đa thức triệt tiêu mạnh.
\end{solution}

\section{Dạng 3: bandwidth và cross-validation}
\begin{exercise}[Bandwidth tham chiếu chuẩn]
Với $X=(1,2,4,5)$, tính bandwidth Gaussian normal-reference.
So sánh hệ số chính xác $(4/3)^{1/5}$ với hệ số làm tròn 1.06.
\end{exercise}
\begin{solution}
\step{1. Tính độ lệch chuẩn mẫu.}
$\overline X=3$, tổng bình phương độ lệch là 10.
\[
s^2=\frac{10}{4-1}=\frac{10}{3},\qquad s\approx1.82574186.
\]
\step{2. Tính thừa số cỡ mẫu.}
$4^{-1/5}\approx0.75785828$. Do đó
\[
\boxed{h_{\rm NR}=(4/3)^{1/5}\sqrt{10/3}\,4^{-1/5}
\approx1.46559887.}
\]
Dùng 1.06 cho $h\approx1.46667281$. Khác biệt là làm tròn hệ số,
không phải phương pháp khác.
\step{3. Nếu dùng scale bền vững.}
Theo quantile nghịch đảo EDF, $Q_{0.25}=1$, $Q_{0.75}=4$, nên
$\IQR=3$, $\IQR/1.34\approx2.23881>s$.
Vậy $a=\min(s,\IQR/1.34)=s$. Quy tắc hệ số $0.9$ cho
$0.9s4^{-1/5}\approx1.24528823$.
\pitfall{Dùng phương sai plug-in $2.5$ thay $s^2=10/3$ làm thay đổi đáp số.
Phải ghi rõ định nghĩa scale và quy ước quantile.}
\end{solution}

\begin{exercise}[Uniform LOOCV bằng phân số]
Cho $X=(0,1,2)$, kernel Uniform. Tính $\widehat J(h)$ với
$h=3/4,3/2$, rồi chọn trong hai ứng viên.
\end{exercise}
\begin{solution}
\step{1. Tính tích phân bằng độ dài giao.}
Mỗi cửa sổ dài $2h$, cao $1/(2nh)$. Hai cửa sổ cách tâm $d$ giao trên
đoạn dài $(2h-|d|)_+$. Do đó
\[
\int\widehat f_h^2=\frac1{4n^2h^2}\sum_{i,j}(2h-|X_i-X_j|)_+.
\]
Trong chín cặp có thứ tự: ba cặp khoảng cách 0, bốn cặp khoảng cách 1,
hai cặp khoảng cách 2.

\step{2. Bandwidth $h=3/4$.}
Tổng độ dài giao là $3(3/2)+4(1/2)=13/2$.
Tích phân bình phương bằng $(13/2)/(36\cdot9/16)=26/81$.
Không có quan sát khác trong khoảng cách $h$ của mỗi điểm, nên tổng
LOO bằng 0:
\[
\boxed{\widehat J(3/4)=26/81\approx0.32098765.}
\]

\step{3. Bandwidth $h=3/2$.}
Tổng độ dài giao là $3(3)+4(2)+2(1)=19$.
Tích phân bằng $19/(36\cdot9/4)=19/81$.
Số hàng xóm LOO lần lượt là $(1,2,1)$, nên các mật độ LOO bằng
$(1/6,1/3,1/6)$ và tổng bằng $2/3$.
\[
\boxed{\widehat J(3/2)=\frac{19}{81}-\frac23\frac23
=-\frac{17}{81}\approx-0.20987654.}
\]
Chọn $\boxed{h=3/2}$ trong hai ứng viên vì tiêu chuẩn nhỏ hơn.
\pitfall{LOO đếm khoảng cách $\leq h$, còn phần tích phân bình phương
phụ thuộc khoảng giao khi khoảng cách $<2h$. Các bài này tránh trường
hợp khoảng cách bằng đúng $h$, nên quy ước đầu mút không đổi điểm CV.}
\end{solution}

\begin{exercise}[Gaussian LOOCV bằng tổng đôi]
Cho $X=(0,1)$, Gaussian. Tính CV tại $h=1/2,1$.
\end{exercise}
\begin{solution}
\step{1. Dùng kernel tích chập.}
$C_K(v)=e^{-v^2/4}/(2\sqrt\pi)$. Có hai cặp đường chéo và hai cặp
ngoài đường chéo, nên
\[
\int\widehat f_h^2=\frac{2C_K(0)+2C_K(1/h)}{4h}
=\frac{1+e^{-1/(4h^2)}}{4\sqrt\pi h}.
\]
\step{2. Mật độ LOO.}
Sau khi bỏ một điểm chỉ còn một điểm; mỗi giá trị LOO bằng
$\varphi(1/h)/h$. Vì $2/n=1$, số hạng trừ là
$2\varphi(1/h)/h$. Do đó
\[
\boxed{\widehat J(h)=\frac{1+e^{-1/(4h^2)}}{4\sqrt\pi h}
-\frac{2\varphi(1/h)}h.}
\]
\step{3. Thế số.}
\begin{center}
\begin{tabular}{cccc}
\toprule $h$ & $\int\widehat f_h^2$ & Tổng LOO & $\widehat J(h)$\\\midrule
$1/2$ & $0.38587167$ & $0.21596387$ & $0.16990780$\\
$1$ & $0.25089522$ & $0.48394145$ & $-0.23304623$\\\bottomrule
\end{tabular}
\end{center}
Chọn $\boxed{h=1}$ trong hai ứng viên. Với mẫu rất nhỏ, quyết định chỉ
minh họa phép tính; không cho biết bandwidth tối ưu toàn bộ miền dương.
\end{solution}

\section{Dạng 4: kiểm định KS từng bước}
\begin{exercise}[{Mẫu trên $[0,1]$}]\label{ex:ksuniform}
Cho $X=(0.1,0.4,0.6,0.8)$, kiểm định $H_0:X\sim U[0,1]$.
Tính $D_n^+$, $D_n^-$, $D_n$. Dùng DKW mức 5\% để quyết định.
\end{exercise}
\begin{solution}
\step{1. Sắp xếp và tính CDF null.}
Mẫu đã sắp, $F_0(X_{(i)})=X_{(i)}$.
\begin{center}
\begin{tabular}{ccccc}
\toprule $i$ & $u_i$ & $i/4-u_i$ & $u_i-(i-1)/4$\\\midrule
1 & $0.10$ & $0.15$ & $0.10$\\
2 & $0.40$ & $0.10$ & $0.15$\\
3 & $0.60$ & $0.15$ & $0.10$\\
4 & $0.80$ & $0.20$ & $0.05$\\\bottomrule
\end{tabular}
\end{center}
\step{2. Lấy cực đại của hai cột sai khác.}
\[
\boxed{D_4^+=0.20,\quad D_4^-=0.15,\quad D_4=0.20.}
\]
\step{3. So với ngưỡng bảo thủ hữu hạn mẫu.}
$\epsilon_4=\sqrt{\log40/8}\approx0.67905$.
$0.20<0.67905$, nên không bác bỏ bằng kiểm định DKW này.
Ngưỡng lớn phản ánh mẫu chỉ có bốn quan sát; không thể kết luận phân
phối đều đúng vì không bác bỏ.

\step{4. Kiểm tra thêm CvM và AD trên cùng mẫu.}
Các tâm so sánh CvM lần lượt là
\[
(0.125,0.375,0.625,0.875).
\]
\[
W_4^2=\frac1{48}+(-.025)^2+(.025)^2+(-.025)^2+(-.075)^2
=\boxed{\frac{17}{600}\approx0.02833333}.
\]
AD dùng công thức log ở Chương 1 cho
$\boxed{A_4^2\approx0.20438803}$.
Các số này là thống kê; cần phân phối null tương ứng để đưa ra quyết định
CvM hoặc AD, không so chúng với ngưỡng KS.
\end{solution}

\begin{exercise}[Null mũ được chỉ định đầy đủ]
Cho $X=(0.2,0.5,1,1.5)$, kiểm định với $F_0(x)=1-e^{-x}$,
$x\geq0$, tức tốc độ mũ cố định bằng 1. Tính KS.
\end{exercise}
\begin{solution}
\step{1. Tính $u_i$.}
\begin{center}
\begin{tabular}{ccccc}
\toprule $i$ & $X_{(i)}$ & $u_i$ & $i/4-u_i$ & $u_i-(i-1)/4$\\\midrule
1 & $0.2$ & $.18126925$ & $.06873075$ & $.18126925$\\
2 & $0.5$ & $.39346934$ & $.10653066$ & $.14346934$\\
3 & $1.0$ & $.63212056$ & $.11787944$ & $.13212056$\\
4 & $1.5$ & $.77686984$ & $.22313016$ & $.02686984$\\\bottomrule
\end{tabular}
\end{center}
\step{2. Cực đại hai phía.}
\[
\boxed{D_4^+=e^{-1.5}\approx0.22313016,\quad
D_4^-=1-e^{-0.2}\approx0.18126925,\quad D_4\approx0.22313016.}
\]
Nếu dùng DKW 5\%, không bác bỏ vì nhỏ hơn $0.67905$.
Không thay tốc độ bằng $1/\overline X$ trong bài này: đề đã chỉ định
tốc độ 1. Nếu fit tốc độ từ mẫu, cần phân phối null được hiệu chỉnh.
\end{solution}

\begin{exercise}[Dữ liệu trùng và giới hạn trái]
Cho $X=(0.2,0.2,0.7,0.9)$, null đều trên $[0,1]$.
Tính KS và giải thích cách xử lý bước nhảy cao $2/4$.
\end{exercise}
\begin{solution}
\step{1. Viết hai cột với tất cả thứ tự thống kê.}
Các $u_i=(.2,.2,.7,.9)$.
\[
\left(\frac i4-u_i\right)=(.05,.30,.05,.10),
\quad
\left(u_i-\frac{i-1}4\right)=(.20,-.05,.20,.15).
\]
Các giá trị ở thứ hạng giữa một nhóm trùng không tương ứng EDF thực
ở một khoảng riêng, nhưng cực đại đúng vẫn được giữ: dùng thứ hạng
cuối nhóm cho phía phải, thứ hạng đầu nhóm cho giới hạn trái.
\step{2. Kiểm tra tại giá trị khác nhau.}
Tại $x=.2$, $\Fhat(.2)=.5$, $\Fhat(.2-)=0$.
Sai khác bên phải là $.5-.2=.3$, bên trái là $.2-0=.2$.
Tại $.7$, hai sai khác là $.75-.7=.05$ và $.7-.5=.2$;
tại $.9$, là $.1$ và $.15$. Vậy
\[
\boxed{D_4^+=.3,\quad D_4^-=.2,\quad D_4=.3.}
\]
Nếu dữ liệu trùng do làm tròn hoặc do phân phối rời rạc, mô hình quan
sát cần được xét lại khi hiệu chỉnh mức kiểm định. Trong một mô hình
liên tục chính xác, xác suất có các giá trị trùng đúng bằng nhau là 0.
Phép tính thống kê vẫn xác định; diễn giải null phải phù hợp dữ liệu.
\end{solution}

\section{Bốn bài chứng minh củng cố}
\begin{exercise}[EDF, covariance và DKW]
Chứng minh covariance EDF khi $x<y$ bằng số đếm; tìm cỡ mẫu để băng
DKW 95\% có nửa độ rộng không quá $0.05$.
\end{exercise}
\begin{solution}
\step{1. Chia ba vùng loại trừ nhau.}
Đặt $A_i=\1_{X_i\leq x}$, $B_i=\1_{x<X_i\leq y}$.
Với $p=F(x)$, $q=F(y)$, $\E A_i=p$, $\E B_i=q-p$,
và $A_iB_i=0$. Vì thế
$\Cov(A_i,B_i)=-p(q-p)$.
\step{2. Dùng $\1_{X_i\leq y}=A_i+B_i$.}
\[
\Cov(A_i,A_i+B_i)=p(1-p)-p(q-p)=p(1-q).
\]
Giữa quan sát khác nhau, covariance bằng 0. Chia tổng covariance cho
$n^2$ cho $\boxed{p(1-q)/n}$.
\step{3. Giải điều kiện độ rộng.}
\[
\sqrt{\frac{\log40}{2n}}\leq.05
\iff n\geq\frac{\log40}{2(.05)^2}\approx737.77589.
\]
Lấy số nguyên lớn hơn hoặc bằng: $\boxed{n\geq738}$.
Không làm tròn xuống vì cần bảo đảm bất đẳng thức.
\end{solution}

\begin{exercise}[Histogram: MISE chính xác cho mật độ tuyến tính]
Cho $f(x)=2x$ trên $[0,1]$, $n=96$, $m=4$.
Tính squared bias tích phân và phương sai tích phân chính xác, rồi
so với AMISE.
\end{exercise}
\begin{solution}
\step{1. Xác suất từng ô.}
$h=1/4$ và
\[
p_j=\int_{(j-1)/4}^{j/4}2x\,dx
=\frac{j^2-(j-1)^2}{16}=\frac{2j-1}{16}.
\]
Vì thế $p=(1,3,5,7)/16$; kỳ vọng histogram ở ô $j$ là $p_j/h=(2j-1)/4$,
bằng $2c_j$. Bias tại $x$ là $2(c_j-x)$, chính xác vì mật độ tuyến tính.
\step{2. Squared bias tích phân.}
\[
\int b^2=\sum_{j=1}^4\int_{-h/2}^{h/2}4t^2dt
=4\frac{4h^3}{12}=\boxed{\frac1{48}}.
\]
\step{3. Phương sai tích phân.}
$\sum_jp_j^2=(1+9+25+49)/256=21/64$.
\[
\int\Var=\frac1{96/4}\left(1-\frac{21}{64}\right)
=\boxed{\frac{43}{1536}}.
\]
Tổng $\MISE=1/48+43/1536=\boxed{25/512\approx0.04882813}$.
\step{4. So với AMISE.}
$R(f')=4$ nên $h_*=(6/(4\cdot96))^{1/3}=1/4$.
AMISE tại đó là $4h^2/12+1/(nh)=1/48+1/24=\boxed{1/16}$.
AMISE khác MISE vì bỏ số hạng âm bậc $1/n$ trong variance tích phân.
\end{solution}

\begin{exercise}[AMISE và độ nhạy bandwidth]
Với Epanechnikov, $R(f'')=1$, $n=480$, tìm $h_*$, AMISE tối ưu và
AMISE tại $h_*/2$, $2h_*$; chứng minh tỷ số tổng quát.
\end{exercise}
\begin{solution}
\step{1. Thế các hằng số kernel.}
\[
\AMISE(h)=\frac{(1/5)^2}{4}h^4+\frac{3/5}{480h}
=\frac{h^4}{100}+\frac1{800h}.
\]
\step{2. Đạo hàm và nghiệm.}
$4h^3/100-1/(800h^2)=0\iff h^5=1/32$, nên
$\boxed{h_*=1/2}$. Khi đó squared bias bằng $1/1600$,
variance bằng $1/400$, tổng $\boxed{1/320}$.
\step{3. Hai ứng viên lệch tối ưu.}
\[
\AMISE(1/4)=\frac1{25600}+\frac1{200}
=\boxed{\frac{129}{25600}},\quad
\AMISE(1)=\frac1{100}+\frac1{800}=\boxed{\frac9{800}}.
\]
\step{4. Tỷ số tổng quát.}
Viết AMISE là $Ah^4+B/(nh)$. Tại nghiệm,
$B/(nh_*)=4Ah_*^4$. Vì vậy
\[
\boxed{\frac{\AMISE(ch_*)}{\AMISE(h_*)}
=\frac{c^4Ah_*^4+4Ah_*^4/c}{5Ah_*^4}
=\frac{c^4+4/c}{5}.}
\]
Đổi bandwidth nửa lần hoặc hai lần không gây thay đổi đối xứng về rủi ro,
vì bias bình phương tăng theo lũy thừa bốn.
\end{solution}

\begin{exercise}[Contrast, Pythagoras và tái tạo tuyến tính]
Chứng minh ba sự triệt tiêu sau: contrast có cực tiểu tại hệ số mẫu;
hạng chéo sai số chiếu bằng 0; local linear triệt tiêu thành phần độ dốc.
Nêu nguyên nhân khác nhau của mỗi sự triệt tiêu.
\end{exercise}
\begin{solution}
\step{1. Contrast.}
Với $g=\sum_{j\leq D}b_j\phi_j$,
\[
\gamma_n(g)=\sum_j(b_j^2-2b_j\widehat a_j)
=\sum_j(b_j-\widehat a_j)^2-\sum_j\widehat a_j^2.
\]
Tổng bình phương cực tiểu bằng 0 khi $b_j=\widehat a_j$.
Đó là hoàn thành bình phương, không phải một giả thiết độc lập.
\step{2. Pythagoras.}
$f-f_D$ có tích vô hướng 0 với mọi $\phi_j$, $j\leq D$;
$f_D-\widehat f_D$ là tổ hợp của các hàm này. Do tuyến tính,
$\inner{f-f_D}{f_D-\widehat f_D}=0$. Khi khai triển bình phương chuẩn,
hạng chéo mất do \emph{trực giao trong không gian hàm}.
\step{3. Local linear.}
Với $w_i=k_i(S_2-S_1d_i)/\Delta$,
\[
\sum_iw_id_i=\frac{S_2\sum_i k_id_i-S_1\sum_i k_id_i^2}{\Delta}
=\frac{S_2S_1-S_1S_2}{\Delta}=0.
\]
Taylor $r(X_i)=r(x)+r'(x)d_i+O(h^2)$;
$r'(x)\sum_iw_id_i=0$. Đây là \emph{đẳng thức moment của trọng số},
đúng với thiết kế cụ thể nếu $\Delta>0$.
\step{4. Kết luận.}
Ba thao tác cùng loại bỏ hạng nhưng lý do khác nhau: tối ưu đại số,
trực giao hàm, và ràng buộc moment. Không viết chung một câu
``do các biến độc lập'' cho cả ba.
\end{solution}

\section{Cách tự chấm bài}
Một lời giải đủ cần có định nghĩa estimator, bảng đối số hoặc ô, mẫu số
đúng, miền áp dụng, phép cộng/tích phân và kết luận.
Với đạo hàm, kiểm tra điểm nối trước khi tính.
Với bandwidth, ghi đây là tối ưu AMSE tại điểm, AMISE toàn cục hay CV
trên tập ứng viên. Với kiểm định, ghi null, thống kê, ngưỡng và quyết
định; tránh biến ``không bác bỏ'' thành ``chứng minh đúng''.
